\hypertarget{c-extensions-python-c-api-interoperability}{%
\chapter{C Extensions \& Python C-API
Interoperability}\label{c-extensions-python-c-api-interoperability}}

\hypertarget{writing-a-pure-c-extension-module}{%
\subsection{24.1 Writing a Pure C Extension
Module}\label{writing-a-pure-c-extension-module}}

Python allows writing modules directly in C or C++ to access low-level
OS APIs, optimize performance-critical paths, or link with native
libraries. A C extension is compiled into a shared library (a
\passthrough{\lstinline!.so!} file on Unix, or a
\passthrough{\lstinline!.pyd!} file on Windows) that Python can load
dynamically at runtime using \passthrough{\lstinline!import!}.

\hypertarget{detailed-c-extension-with-error-handling-and-keywords}{%
\subsubsection{1. Detailed C-Extension with Error Handling and
Keywords}\label{detailed-c-extension-with-error-handling-and-keywords}}

A robust C extension must handle keyword arguments, perform error
checks, raise Python exceptions on failure, and manage resources.

Below is an implementation of a custom C module exposing a safe division
function:

\begin{lstlisting}[language=C]
#define PY_SSIZE_T_CLEAN
#include <Python.h>

/* 1. Core C Function Implementation with Keywords & Exception Handling */
static PyObject* 
custom_safe_divide(PyObject* self, PyObject* args, PyObject* keywds) {
    double numerator;
    double denominator;
    
    /* Argument keyword names array */
    static char *kwlist[] = {"numerator", "denominator", NULL};
    
    /* Parse Python argument tuple & dict into native double values */
    if (!PyArg_ParseTupleAndKeywords(args, keywds, "dd", kwlist, &numerator, &denominator)) {
        return NULL; /* Returns NULL to signal an exception occurred during parsing */
    }
    
    /* Check for division by zero and raise a Python ZeroDivisionError */
    if (denominator == 0.0) {
        PyErr_SetString(PyExc_ZeroDivisionError, "Denominator cannot be zero.");
        return NULL; /* Return NULL to propagate the raised exception */
    }
    
    double result = numerator / denominator;
    
    /* Convert native C double back to a Python PyFloatObject */
    return PyFloat_FromDouble(result);
}

/* 2. Methods Table Registration */
static PyMethodDef CustomMethods[] = {
    {
        "safe_divide", 
        (PyCFunction)(void(*)(void))custom_safe_divide, 
        METH_VARARGS | METH_KEYWORDS, 
        "Divides numerator by denominator, returning float safely."
    },
    {NULL, NULL, 0, NULL} /* Sentinel marker to end array */
};

/* 3. Module Definition Structure */
static struct PyModuleDef custommodule = {
    PyModuleDef_HEAD_INIT,
    "custom_c_math",            /* Module Name */
    "A custom high-performance C extension module.", /* Docstring */
    -1,                         /* Module state size (-1 = global state) */
    CustomMethods
};

/* 4. Initialization Hook called upon 'import custom_c_math' */
PyMODINIT_FUNC 
PyInit_custom_c_math(void) {
    return PyModule_Create(&custommodule);
}
\end{lstlisting}

\begin{itemize}
\tightlist
\item
  \textbf{\passthrough{\lstinline!METH\_VARARGS | METH\_KEYWORDS!}}:
  Tells the interpreter that this function accepts both positional
  arguments (parsed into \passthrough{\lstinline!args!} as a tuple) and
  keyword arguments (parsed into \passthrough{\lstinline!keywds!} as a
  dictionary).
\item
  \textbf{\passthrough{\lstinline!PyArg\_ParseTupleAndKeywords!}}:
  Resolves both types of inputs, mapping them into the C double
  variables based on the format string \passthrough{\lstinline!"dd"!}
  (double, double) and the keyword list
  \passthrough{\lstinline!kwlist!}.
\item
  \textbf{Exception Signaling}: Returning a
  \passthrough{\lstinline!NULL!} pointer tells CPython's execution frame
  that an exception has been set inside the thread state. The
  interpreter pauses bytecode execution and executes its exception
  handler path.
\end{itemize}

\begin{center}\rule{0.5\linewidth}{0.5pt}\end{center}

\hypertarget{reference-counting-and-ownership-rules-in-c-api}{%
\subsection{24.2 Reference Counting and Ownership Rules in
C-API}\label{reference-counting-and-ownership-rules-in-c-api}}

When interacting with the Python C-API, developers must manually manage
reference counting. Standard Python code delegates reference updates to
the compiler/interpreter, but in C, a single missing decref causes a
memory leak, and an extra decref causes a crash (segmentation fault).

\hypertarget{reference-ownership-categories}{%
\subsubsection{1. Reference Ownership
Categories}\label{reference-ownership-categories}}

Every pointer returning from a C-API function belongs to one of three
reference categories:

\hypertarget{new-references}{%
\paragraph{New References}\label{new-references}}

The C function receives absolute ownership of the reference. It must
decrement the reference count when finished, or pass ownership back to
CPython (e.g., by returning the object from the function).

\begin{itemize}
\tightlist
\item
  \emph{Examples}: \passthrough{\lstinline!PyLong\_FromLong()!},
  \passthrough{\lstinline!PyList\_New()!},
  \passthrough{\lstinline!PyObject\_Call()!}.
\end{itemize}

\hypertarget{borrowed-references}{%
\paragraph{Borrowed References}\label{borrowed-references}}

The C function receives a pointer without ownership. It does not own the
reference and must not decrement it unless it explicitly calls
\passthrough{\lstinline!Py\_INCREF()!} to claim ownership. If the owner
of the object frees it, a borrowed pointer becomes dangling.

\begin{itemize}
\tightlist
\item
  \emph{Examples}: \passthrough{\lstinline!PyTuple\_GetItem()!},
  \passthrough{\lstinline!PyList\_GetItem()!}.
\end{itemize}

\hypertarget{stolen-references}{%
\paragraph{Stolen References}\label{stolen-references}}

Some C-API functions take ownership of references passed to them. The
caller no longer owns the reference and must not decrement it, as the
receiver will decrement it during its own deallocation pass.

\begin{itemize}
\tightlist
\item
  \emph{Examples}: \passthrough{\lstinline!PyTuple\_SetItem()!},
  \passthrough{\lstinline!PyList\_SetItem()!}.
\end{itemize}

\begin{lstlisting}[language=C]
void add_item_to_tuple_example(void) {
    PyObject* my_tuple = PyTuple_New(1);              /* Returns New Reference */
    PyObject* my_val = PyLong_FromLong(42);           /* Returns New Reference */
    
    /* PyTuple_SetItem steals the reference to 'my_val' */
    PyTuple_SetItem(my_tuple, 0, my_val);
    
    /* 
       Do NOT call Py_DECREF(my_val). 'my_tuple' now owns it.
       If we decref'd my_val, freeing my_tuple would trigger a double-free crash.
    */
    
    Py_DECREF(my_tuple); /* Safely deallocates tuple and my_val */
}
\end{lstlisting}

\hypertarget{safe-reference-macros}{%
\subsubsection{2. Safe Reference Macros}\label{safe-reference-macros}}

CPython provides macros to handle reference updates safely: *
\passthrough{\lstinline!Py\_INCREF(op)!} /
\passthrough{\lstinline!Py\_DECREF(op)!}: Non-null safe updates. Passing
a \passthrough{\lstinline!NULL!} pointer triggers a crash. *
\passthrough{\lstinline!Py\_XINCREF(op)!} /
\passthrough{\lstinline!Py\_XDECREF(op)!}: Null-safe updates. They check
if \passthrough{\lstinline!op!} is \passthrough{\lstinline!NULL!} and do
nothing if it is, protecting code in error-handling block cleanups:

\begin{lstlisting}[language=C]
void error_handling_cleanup_example(PyObject *a, PyObject *b) {
    // If a or b failed allocation and were set to NULL, XDECREF handles it safely
    Py_XDECREF(a);
    Py_XDECREF(b);
}
\end{lstlisting}

\begin{center}\rule{0.5\linewidth}{0.5pt}\end{center}

\hypertarget{c-api-stable-abi-pep-384}{%
\subsection{24.3 C-API Stable ABI (PEP
384)}\label{c-api-stable-abi-pep-384}}

Standard C extensions are compiled against a specific Python version's
headers (e.g., Python 3.10). They depend on concrete structure offsets
(such as \passthrough{\lstinline!sizeof(PyObject)!} or offsets of
\passthrough{\lstinline!ob\_type!}). When Python updates its internal
layouts (e.g., changing the type struct in 3.11), these compiled
binaries break, requiring recompilation for the new Python version.

To resolve this compilation dependency, PEP 384 introduced the
\textbf{Stable ABI} (Application Binary Interface) / Limited API.

\begin{lstlisting}
+---------------------------------------+
|        Standard C Extension           |
|  Directly accesses structures (e.g.   | ---> Fast, but must compile
|  op->ob_refcnt, op->ob_type)          |      for every Python version.
+---------------------------------------+

+---------------------------------------+
|          Stable ABI Extension         |
|  Uses opaque pointers and accessors   | ---> Compile once. Runs on
|  (e.g., Py_REFCNT(op), PyType_GetSlot) |      Python 3.5, 3.6, ... 3.13+
+---------------------------------------+
\end{lstlisting}

\hypertarget{opaque-pointers-and-opaque-structures}{%
\subsubsection{1. Opaque Pointers and Opaque
Structures}\label{opaque-pointers-and-opaque-structures}}

When compiling under the Stable ABI, structure definitions (like
\passthrough{\lstinline!PyObject!} or
\passthrough{\lstinline!PyTypeObject!}) are treated as \textbf{opaque
structures}. You cannot compile code that accesses their internal fields
directly:

\begin{lstlisting}[language=C]
/* Standard API (will not compile under Py_LIMITED_API) */
PyTypeObject* type = op->ob_type;

/* Stable ABI equivalent */
PyTypeObject* type = Py_TYPE(op);
\end{lstlisting}

\hypertarget{defining-the-limited-api}{%
\subsubsection{2. Defining the Limited
API}\label{defining-the-limited-api}}

To restrict compilation to the Stable ABI, define
\passthrough{\lstinline!Py\_LIMITED\_API!} before importing
\passthrough{\lstinline!<Python.h>!} or set it as a compiler flag:

\begin{lstlisting}[language=C]
#define Py_LIMITED_API 0x030A0000  /* Targets Python 3.10 and later */
#define PY_SSIZE_T_CLEAN
#include <Python.h>
\end{lstlisting}

This guarantees that the resulting compiled binary links only with
stable symbols exported by the dynamic library
(\passthrough{\lstinline!libpython3.so!} or
\passthrough{\lstinline!python3.dll!}), allowing the extension to load
and run on any subsequent Python version without recompilation.

\begin{center}\rule{0.5\linewidth}{0.5pt}\end{center}
